\section{Exact finite core and paper-level identities}\label{sec:proof}
\subsection{Lean-verified finite algebra}
The pinned Lean~4 / Mathlib project under
\texttt{formalisation/} contains exact finite theorems,
not asymptotic estimates. The twenty-six independently
compared selections are listed with precise names and
hypotheses in \texttt{formalisation/README.md};
the comparison includes:
\texttt{sexticExp\_mul\_sixth\_pow},
\texttt{sexticExp\_pow\_five},
\texttt{rawGauss\_modulus\_sq},
\texttt{shiftedGauss\_e4}, and
\texttt{residueMatVec\_frozen\_orientation}.
The checked source baseline is commit
\texttt{8b050b06d1195491b17b92e68b2d92d9d2478cb0}
and the independent Comparator/NanoDa success run is
\href{https://github.com/jfairfaxball-348/PRIMES/actions/runs/40569785577}{40569785577};
the corresponding full scoped cloud workflow is
\href{https://github.com/jfairfaxball-348/PRIMES/actions/runs/40569785692}{40569785692}.
The code also reports \(1764/1764\) agreement in its
finite exponent-faithfulness fixture.
% Source: PROGRAM_STATUS.md RS-S040 checkpoint; formalisation/README.md (main 30965467).
The formal residue convention maps
\(\omega\mapsto r=-a/b\) in \(\mathrm{ZMod}(p)\),
with explicit prime/root/divisibility hypotheses;
it is \emph{not} a Mathlib construction of
the quotient \(\OO/(\sigma)\).
The frozen CAND-002 statements are defined as
Lean propositions, with \emph{no proofs} of
their asymptotic parts. Palomar registration is
not asserted.

\subsection{Exact shifted row and signed Gram identities}
Put \(G_\sigma=\gtilde(\sigma)\),
\(q_\kappa(\sigma)=\overline{G_\sigma}\chi_\sigma(\kappa)\),
and \(U_\kappa=\{\sigma\in\PP(X):\chi_\sigma(\kappa)\ne0\}\);
let \(D_\kappa=|U_\kappa|\).
For the positive-sign rational additive character,
\[
 g_\sigma(u)=\sum_{a\bmod\sigma}
    \chi_\sigma(a)e(\operatorname{res}_\sigma(ua)/\Nm\sigma).
\]

\begin{proposition}[Finite row identity, exact]\label{prop:row}
For every indexed row \(\rho\), with the zero convention,
\[
 (\rX w_\kappa)_\rho=
 \overline{\sum_{\sigma\in\PP(X)}
              \frac{g_\sigma(\rho\kappa)}{\sqrt{\Nm\sigma}}}.
\]
\end{proposition}
\begin{proof}
For \(\sigma\nmid \rho\kappa\), substitute
\(a\mapsto(\rho\kappa)^{-1}a\) in the Gauss sum:
\(g_\sigma(\rho\kappa)
=\overline{\chi_\sigma(\rho\kappa)}g_\sigma(1)\).
Conjugating and multiplying symbols gives exactly the
\(\sigma\)-summand of the matrix product. If
\(\sigma\mid\rho\kappa\), the shifted sum vanishes by
nontrivial multiplicative-character orthogonality,
and so does the corresponding matrix term.
This includes the zero diagonal.
\end{proof}

Define a Hermitian prime-correlation kernel and real signed excess
\begin{align*}
 K_X(\sigma,\tau)
   &=\sum_{\rho\in\PP(X)}
          \chi_\sigma(\rho)\overline{\chi_\tau(\rho)},\\
 O_\kappa
   &=\sum_{\substack{\sigma,\tau\in U_\kappa\\\sigma\ne\tau}}
      q_\kappa(\sigma)\overline{q_\kappa(\tau)}
         K_X(\sigma,\tau).
\end{align*}
\begin{proposition}[Exact zero-diagonal Gram expansion]\label{prop:gram}
For \(D,D_\kappa>0\),
\[
 \|\rX w_\kappa\|^2=D_\kappa(D-1)+O_\kappa,\qquad
 r_\kappa^2=1-\frac1D+\frac{O_\kappa}{DD_\kappa}.
\]
Consequently, wherever \(r_1^2\ne1\),
\[
 s_\kappa=\Nm_6(\kappa)^{1/6}
 \frac{O_\kappa/D_\kappa-1}{O_1/D-1}.
\]
\end{proposition}
\begin{proof}
Expand the squared row sums and interchange finite sums.
For a diagonal column pair \(\sigma=\tau\),
\(|q_\kappa(\sigma)|^2=1\), while the
inner row sum has \(D-1\) unit-modulus terms
and one zero term at \(\rho=\sigma\).
The off-diagonal pairs give \(O_\kappa\).
Conjugate-symmetric pairing makes \(O_\kappa\) real;
normalising gives the formula.
\end{proof}
These results follow RS-S026 and do not assume an
unproved first-moment estimate. Their finite row
counterpart is independently Lean-checked;
the separate Gram expansion is proved here on paper.

\subsection{Reciprocity kernel, sieve and logical limitation}
For distinct columns away from exceptional local primes,
RS-S027 gives an exact reciprocity transform
\[
 K_X(\sigma,\tau)
  =\sum_{\rho\in\PP(X)}
       h_{\sigma,\tau}(\rho)
       \left(\frac{\sigma\tau^5}{\rho}\right)_6,
\]
where \(h_{\sigma,\tau}\) is a \emph{quadratic
local factor modulo \(4\)} determined by the
primary representatives. This is a rearrangement,
not cancellation. The conductor grows with
\(\sigma,\tau\), so the fixed-character prime-ideal
theorem does not establish the required
uniform signed estimate.

The ideal-prime projection in RS-S028 applies
a von Mangoldt weight and separates inert ideals,
prime powers and local exclusions. An exact ideal
Vaughan convolution decomposes the weighted
sum into Type I and Type II pieces.
A separate Dirichlet-series countermodel in
that audit shows that knowledge of a pole
\emph{residue alone} does not determine
a prime-sieved coefficient. These are proofs
of reductions and obstructions, not a proof of
a sextic coefficient asymptotic.

Likewise, the abstract scalar countermodel in
\texttt{RS\_CAND002\_LOGICAL\_BRIDGE\_AUDIT.md}
demonstrates that compact-positive-aspect and
large-aspect hypotheses need not specify the
fixed-class boundary. It is \emph{not}
a counterexample to the arithmetic conjecture.
Detailed mod-\(4\) local signs, trace-to-rational
normalisation corrections, and the ideal
Vaughan terms remain in the referenced
paper-level derivations; no unproved transfer
has been promoted here.
